% ============================================================
%  RouteForge -- Academic Project Report
%  Compile with: pdflatex -> pdflatex
%  or: latexmk -pdf routeforge_report.tex
% ============================================================

\documentclass[12pt, a4paper]{article}

% -- Packages ------------------------------------------------
\usepackage[utf8]{inputenc}
\usepackage[T1]{fontenc}
\usepackage{lmodern}
\usepackage[british]{babel}
\usepackage[margin=2.5cm]{geometry}
\usepackage{setspace}
\usepackage{parskip}
\usepackage{titlesec}
\usepackage{fancyhdr}
\usepackage{float}
\usepackage{booktabs}
\usepackage{tabularx}
\usepackage{array}
\usepackage{amsmath}
\usepackage{amssymb}
\usepackage{algorithm}
\usepackage{algpseudocode}
\usepackage{listings}
\usepackage{xcolor}
\usepackage{hyperref}
\usepackage{caption}
\usepackage{enumitem}
\usepackage{microtype}
\usepackage{tocloft}

% -- Hyperref ------------------------------------------------
\hypersetup{
    colorlinks  = true,
    linkcolor   = blue!60!black,
    citecolor   = green!50!black,
    urlcolor    = blue!70!black,
    pdftitle    = {RouteForge: A Graph Algorithm Visualisation System},
    pdfsubject  = {Data Structures and Algorithms},
}
% -- Spacing -------------------------------------------------
\onehalfspacing

% -- Section formatting --------------------------------------
\titleformat{\section}{\large\bfseries}{\thesection}{1em}{}[\titlerule]
\titleformat{\subsection}{\normalsize\bfseries}{\thesubsection}{1em}{}
\titleformat{\subsubsection}{\normalsize\itshape}{\thesubsubsection}{1em}{}

% -- Header / footer -----------------------------------------
\pagestyle{fancy}
\fancyhf{}
\fancyhead[L]{\small RouteForge}
\fancyhead[R]{\small Data Structures \& Algorithms}
\fancyfoot[C]{\thepage}
\renewcommand{\headrulewidth}{0.4pt}

% -- Code style ----------------------------------------------
\definecolor{codebg}{RGB}{245,247,250}
\definecolor{codeframe}{RGB}{200,210,220}
\definecolor{kwcolor}{RGB}{0,95,190}
\definecolor{strcolor}{RGB}{160,30,30}
\definecolor{cmtcolor}{RGB}{100,115,130}
\definecolor{numcolor}{RGB}{100,60,160}

\lstdefinestyle{python}{
    language         = Python,
    backgroundcolor  = \color{codebg},
    frame            = single,
    rulecolor        = \color{codeframe},
    basicstyle       = \ttfamily\footnotesize,
    keywordstyle     = \color{kwcolor}\bfseries,
    stringstyle      = \color{strcolor},
    commentstyle     = \color{cmtcolor}\itshape,
    numberstyle      = \tiny\color{numcolor},
    numbers          = left,
    numbersep        = 8pt,
    stepnumber       = 1,
    tabsize          = 4,
    showstringspaces = false,
    breaklines       = true,
    breakatwhitespace= true,
    captionpos       = b,
    aboveskip        = 8pt,
    belowskip        = 4pt,
    xleftmargin      = 14pt,
    framexleftmargin = 14pt,
}
\lstset{style=python}

% -- Algorithmic ---------------------------------------------
\algrenewcommand\algorithmicprocedure{\textbf{procedure}}
\algrenewcommand\algorithmicfunction{\textbf{function}}

% -- Captions ------------------------------------------------
\captionsetup{font=small, labelfont=bf, skip=6pt}

% -- Column types --------------------------------------------
\newcolumntype{L}[1]{>{\raggedright\arraybackslash}p{#1}}
\newcolumntype{C}[1]{>{\centering\arraybackslash}p{#1}}
\newcolumntype{R}[1]{>{\raggedleft\arraybackslash}p{#1}}

% ============================================================
\begin{document}
% ============================================================

% -- Title page ----------------------------------------------
\begin{titlepage}
    \centering
    \vspace*{2cm}

    {\Large\textsc{Year 2 -- Data Structures \& Algorithms}\par}
    \vspace{0.4cm}
    {\large\textsc{Project Report}\par}

    \vspace{1.5cm}
    \rule{\linewidth}{1pt}\\[0.6em]
    {\Huge\bfseries RouteForge\par}
    \vspace{0.3em}
    {\Large A Graph Algorithm Visualisation System\par}
    \rule{\linewidth}{1pt}

    \vspace{2cm}

    \begin{tabular}{rl}
        \textbf{Module:}     & Data Structures \& Algorithms \\[4pt]
        \textbf{Language:}   & Python 3.11+                  \\[4pt]
        \textbf{Framework:}  & pygame 2.x                    \\[4pt]
    \end{tabular}

    \vspace{1.8cm}

    {\large\textbf{Authors}\par}
    \vspace{0.4cm}
    \begin{tabular}{c}
        Blevis Allushi (AI) \\[3pt]
        Kristian Seraj (AI)\\[3pt]
        Renato Zotaj   (SWE)\\[3pt]
        Leandra Latifi (AI)\\
    \end{tabular}

    \vfill
\end{titlepage}

% -- Table of contents ---------------------------------------
\tableofcontents
\newpage

% ============================================================
\section{Introduction}
% ============================================================

Graph theory underpins a broad class of real-world problems -- from network
routing and geographic pathfinding to social-network analysis and dependency
resolution. Understanding how fundamental graph algorithms operate is a core
competency in computer science education; yet the abstract mechanics of
queues, priority queues, and predecessor maps are often difficult to grasp
from code or pseudocode alone.

RouteForge addresses this gap by pairing correct, modular algorithm
implementations with an interactive step-by-step visualiser. The system is
designed around three principles:

\begin{enumerate}[label=(\roman*)]
    \item \textbf{Separation of concerns} -- the graph data structure, the
          algorithms, the CLI, and the visualiser are entirely independent
          modules.
    \item \textbf{Correctness first} -- each algorithm is implemented
          faithfully to its canonical form; the visualiser replicates the
          algorithm logic internally without modifying the production code.
    \item \textbf{Interactivity} -- the user can construct arbitrary graphs,
          generate random instances, and animate algorithms at configurable
          speeds, stepping forwards and backwards through execution.
\end{enumerate}

% ============================================================
\section{System Architecture}
\label{sec:arch}
% ============================================================

\subsection{Module Layout}

The project is organised as a Python package rooted at \texttt{src/}.

\begin{verbatim}
RouteForge/
|-- main.py
`-- src/
    |-- algorithms/
    |   |-- bfs.py
    |   |-- dfs.py
    |   |-- dijkstra.py
    |   |-- helpers.py
    |   |-- kruskal.py
    |   `-- prim.py
    |-- core/
    |   |-- graph.py
    |   `-- validators.py
    |-- data/
    |-- ui/
    |   `-- menu.py
    |-- utils/
    |   `-- graph_generator.py
    `-- visualization/
        `-- pygame_view.py
\end{verbatim}

\subsection{Module Dependencies}

The dependency structure is intentionally acyclic and layered:

\begin{itemize}
    \item \texttt{src/core/} has no internal dependencies; it defines the
          foundation types used by every other layer.
    \item \texttt{src/algorithms/} depends only on \texttt{src/core/}.
    \item \texttt{src/utils/} depends only on \texttt{src/core/}.
    \item \texttt{src/visualization/} depends on \texttt{src/core/} only.
    \item \texttt{src/ui/} depends on all other layers and acts as the
          integration point.
    \item \texttt{main.py} depends only on \texttt{src/ui/menu.py}.
\end{itemize}

This layering ensures that the core graph and algorithm modules can be
tested and used independently of the interface layers.

\subsection{Validation Architecture}

Input validation is split across three layers:

\begin{table}[H]
\centering
\caption{Validation responsibility by layer.}
\begin{tabularx}{\linewidth}{@{}L{2.8cm}L{4.2cm}L{6.2cm}@{}}
\toprule
\textbf{Layer} & \textbf{Module} & \textbf{Responsibility} \\
\midrule
Domain rules
  & \texttt{validators.py}
  & Node name format, weight range, NaN/Inf rejection, self-loop prevention \\
Structural
  & \texttt{graph.py}
  & Node existence, endpoint auto-creation, adjacency consistency \\
UI / interaction
  & \texttt{pygame\_view.py}
  & Scroll-wheel guard, click blocking during input, error message display \\
\bottomrule
\end{tabularx}
\end{table}

% ============================================================
\section{Core Data Structures}
\label{sec:ds}
% ============================================================

\subsection{Graph Representation}

RouteForge models a weighted undirected graph $G = (V, E, w)$ using an
\emph{adjacency list} stored as a nested Python dictionary:

\[
    \texttt{\_adj} : \text{str} \;\to\; (\text{str} \;\to\; \text{float})
\]

Each outer key is a node name; its value maps neighbour names to edge
weights. Every undirected edge $\{u,v,w\}$ is stored as two directed
entries so that \texttt{\_adj[u][v] = \_adj[v][u] = w}.

The adjacency list was chosen over an adjacency matrix because it uses
$O(|V|+|E|)$ space, supports $O(1)$ edge lookup via nested dict membership,
and accommodates dynamic node/edge addition without reallocation.

\begin{table}[H]
\centering
\caption{Time complexity of core \texttt{Graph} operations.}
\begin{tabular}{@{}lll@{}}
\toprule
\textbf{Operation} & \textbf{Complexity} & \textbf{Notes} \\
\midrule
\texttt{add\_node}    & $O(1)$          & Dict insert \\
\texttt{add\_edge}    & $O(1)$          & Two dict inserts + validation \\
\texttt{remove\_node} & $O(\deg(v))$   & Removes reverse entries \\
\texttt{remove\_edge} & $O(1)$          & Two dict deletes \\
\texttt{has\_node}    & $O(1)$          & Dict membership \\
\texttt{has\_edge}    & $O(1)$          & Nested dict membership \\
\texttt{neighbors}    & $O(\deg(v))$   & Returns list of (neighbour, weight) \\
\texttt{edges}        & $O(|V|+|E|)$   & Deduplication via canonical key \\
Space                 & $O(|V|+|E|)$   & \\
\bottomrule
\end{tabular}
\end{table}

\subsection{Edge Dataclass}

Edges are returned as instances of a frozen, hashable \texttt{Edge}
dataclass, enabling their use in sets and as dictionary keys:

\begin{lstlisting}[caption={Edge dataclass (\texttt{src/core/graph.py}).}]
@dataclass(frozen=True)
class Edge:
    u: str
    v: str
    weight: float
\end{lstlisting}

The \texttt{edges()} method deduplicates using a canonical key
$(\min(u,v),\, \max(u,v))$, ensuring each undirected edge appears exactly
once regardless of insertion order.

\subsection{Disjoint Set (Union-Find)}

Kruskal's algorithm requires a disjoint-set structure to detect cycles.
The \texttt{DisjointSet} class implements union by rank with path
compression, giving an amortised cost of $O(\alpha(n))$ per operation
where $\alpha$ is the inverse Ackermann function -- effectively constant
for all practical inputs.

\begin{lstlisting}[caption={Disjoint-Set (\texttt{src/algorithms/kruskal.py}).}]
class DisjointSet:
    def __init__(self, items):
        self.parent = {item: item for item in items}
        self.rank   = {item: 0    for item in items}

    def find(self, x):
        if self.parent[x] != x:
            self.parent[x] = self.find(self.parent[x])  # path compression
        return self.parent[x]

    def union(self, a, b):
        root_a, root_b = self.find(a), self.find(b)
        if root_a == root_b:
            return False
        if self.rank[root_a] < self.rank[root_b]:
            self.parent[root_a] = root_b
        elif self.rank[root_a] > self.rank[root_b]:
            self.parent[root_b] = root_a
        else:
            self.parent[root_b] = root_a
            self.rank[root_a]  += 1
        return True
\end{lstlisting}

% ============================================================
\section{Input Validation}
\label{sec:validation}
% ============================================================

All domain-level rules live in \texttt{validators.py} and are called by
\texttt{add\_node} and \texttt{add\_edge} before any structural change
occurs. The UI layer also calls validators independently to surface
precise error messages in the status bar without crashing.

\begin{itemize}
    \item \textbf{Node names} -- non-empty, at most 16 characters,
          restricted to \texttt{[a-zA-Z0-9\_-]}.
    \item \textbf{Weights} -- numeric, finite (no NaN or $\pm\infty$),
          within $[0.0,\; 10^6]$.
    \item \textbf{Edges} -- both endpoint names must be valid and
          distinct (no self-loops).
\end{itemize}

All validators raise \texttt{ValueError} with a descriptive message on
failure.

% ============================================================
\section{Algorithms}
\label{sec:algorithms}
% ============================================================

% ------------------------------------------------------------
\subsection{Breadth-First Search}
\label{sec:bfs}
% ------------------------------------------------------------

BFS explores a graph level by level using a FIFO queue, visiting all
neighbours of a node before advancing to their neighbours. Neighbours
are processed in lexicographic order to ensure a deterministic traversal.

\begin{algorithm}[H]
\caption{Breadth-First Search}
\begin{algorithmic}[1]
\Function{BFS}{$G$, $start$}
    \State $visited \gets \{start\}$;\ $queue \gets \langle start \rangle$;\ $order \gets []$
    \While{$queue \neq \emptyset$}
        \State $node \gets$ \Call{Dequeue}{$queue$}
        \State $order.\textsc{Append}(node)$
        \For{each $(nbr, \_)$ in \textsc{SortedNeighbours}($G$, $node$)}
            \If{$nbr \notin visited$}
                \State $visited.\textsc{Add}(nbr)$;\ \Call{Enqueue}{$queue$, $nbr$}
            \EndIf
        \EndFor
    \EndWhile
    \Return $order$
\EndFunction
\end{algorithmic}
\end{algorithm}

\texttt{collections.deque} is used for $O(1)$ \texttt{popleft} operations;
a plain list would give $O(n)$ \texttt{pop(0)}.

\begin{lstlisting}[caption={BFS implementation (\texttt{src/algorithms/bfs.py}).}]
def bfs(graph: Graph, start: str) -> List[str]:
    if not graph.has_node(start):
        raise KeyError(f"Start node '{start}' does not exist.")
    visited = set([start])
    queue   = deque([start])
    order: List[str] = []
    while queue:
        node = queue.popleft()
        order.append(node)
        for neighbor, _ in sorted(graph.neighbors(node), key=lambda x: x[0]):
            if neighbor not in visited:
                visited.add(neighbor)
                queue.append(neighbor)
    return order
\end{lstlisting}

\textbf{Complexity:} $O(|V|+|E|)$ time, $O(|V|)$ space.

% ------------------------------------------------------------
\subsection{Depth-First Search}
\label{sec:dfs}
% ------------------------------------------------------------

DFS explores as far as possible along each branch before backtracking.
An iterative implementation with an explicit stack is used to avoid
Python's recursion depth limit. Neighbours are pushed in reverse-sorted
order so the lexicographically smallest is expanded first, matching
recursive DFS behaviour. The visited check occurs on \textit{pop} rather
than \textit{push} because a node may be pushed multiple times before
first expansion.

\begin{algorithm}[H]
\caption{Depth-First Search (Iterative)}
\begin{algorithmic}[1]
\Function{DFS}{$G$, $start$}
    \State $visited \gets \emptyset$;\ $stack \gets [start]$;\ $order \gets []$
    \While{$stack \neq \emptyset$}
        \State $node \gets stack.\textsc{Pop}()$
        \If{$node \in visited$} \textbf{continue} \EndIf
        \State $visited.\textsc{Add}(node)$;\ $order.\textsc{Append}(node)$
        \For{each $(nbr, \_)$ in \textsc{ReverseSortedNeighbours}($G$, $node$)}
            \If{$nbr \notin visited$}
                \State $stack.\textsc{Push}(nbr)$
            \EndIf
        \EndFor
    \EndWhile
    \Return $order$
\EndFunction
\end{algorithmic}
\end{algorithm}

\begin{lstlisting}[caption={DFS implementation (\texttt{src/algorithms/dfs.py}).}]
def dfs(graph: Graph, start: str) -> List[str]:
    if not graph.has_node(start):
        raise KeyError(f"Start node '{start}' does not exist.")
    visited = set()
    stack   = [start]
    order: List[str] = []
    while stack:
        node = stack.pop()
        if node in visited:
            continue
        visited.add(node)
        order.append(node)
        neighbors = sorted(graph.neighbors(node), key=lambda x: x[0], reverse=True)
        for neighbor, _ in neighbors:
            if neighbor not in visited:
                stack.append(neighbor)
    return order
\end{lstlisting}

\textbf{Complexity:} $O(|V|+|E|)$ time, $O(|V|)$ space.

% ------------------------------------------------------------
\subsection{Dijkstra's Shortest-Path Algorithm}
\label{sec:dijkstra}
% ------------------------------------------------------------

Dijkstra's algorithm computes single-source shortest paths in graphs with
non-negative edge weights. It greedily expands the nearest unvisited node
using a min-priority queue and relaxes outgoing edges. The
\texttt{validate\_weight} call in \texttt{add\_edge} enforces $w \geq 0$
at construction time, guaranteeing the correctness precondition.

\begin{algorithm}[H]
\caption{Dijkstra's Algorithm}
\begin{algorithmic}[1]
\Function{Dijkstra}{$G$, $start$}
    \State $dist[v] \gets \infty$ for all $v$;\ $dist[start] \gets 0$;\ $prev[v] \gets \textsc{None}$
    \State $pq \gets \{(0, start)\}$;\ $visited \gets \emptyset$
    \While{$pq \neq \emptyset$}
        \State $(d, node) \gets \textsc{ExtractMin}(pq)$
        \If{$node \in visited$} \textbf{continue} \EndIf
        \State $visited.\textsc{Add}(node)$
        \For{each $(nbr, w)$ in \textsc{Neighbours}($G$, $node$)}
            \If{$d + w < dist[nbr]$}
                \State $dist[nbr] \gets d + w$;\ $prev[nbr] \gets node$
                \State $pq.\textsc{Insert}(d+w,\ nbr)$
            \EndIf
        \EndFor
    \EndWhile
    \Return $dist$, $prev$
\EndFunction
\end{algorithmic}
\end{algorithm}

Lazy deletion is used: stale heap entries are discarded by the
\textit{visited} check on extraction, avoiding the need for a
decrease-key operation.

\begin{lstlisting}[caption={Dijkstra core (\texttt{src/algorithms/dijkstra.py}).}]
def dijkstra(graph, start):
    distances = {node: float("inf") for node in graph.nodes()}
    previous  = {node: None         for node in graph.nodes()}
    distances[start] = 0.0
    pq      = [(0.0, start)]
    visited = set()
    while pq:
        current_dist, node = heapq.heappop(pq)
        if node in visited:
            continue
        visited.add(node)
        for neighbor, weight in graph.neighbors(node):
            new_dist = current_dist + weight
            if new_dist < distances[neighbor]:
                distances[neighbor] = new_dist
                previous[neighbor]  = node
                heapq.heappush(pq, (new_dist, neighbor))
    return distances, previous
\end{lstlisting}

Path reconstruction is handled by \texttt{reconstruct\_path} in
\texttt{helpers.py}, which walks the predecessor map backwards from the
end node and reverses the result. If the end node is unreachable,
\texttt{None} is returned.

\textbf{Complexity:} $O\bigl((|V|+|E|)\log|V|\bigr)$ time, $O(|V|)$ space.

% ------------------------------------------------------------
\subsection{Minimum Spanning Tree Algorithms}
\label{sec:mst}
% ------------------------------------------------------------

Both MST algorithms are exposed through the CLI and visualised
step-by-step in the pygame interface.

\subsubsection{Prim's Algorithm}

Prim's algorithm grows the MST one edge at a time, repeatedly selecting
the minimum-weight edge crossing the frontier between visited and
unvisited nodes. It requires a connected graph and an optional start node
(defaults to the first node).

\begin{algorithm}[H]
\caption{Prim's MST}
\begin{algorithmic}[1]
\Function{PrimMST}{$G$, $start$}
    \State $visited \gets \{start\}$;\ $mst \gets []$;\ push all edges from $start$ into $pq$
    \While{$pq \neq \emptyset$ \textbf{and} $|visited| < |V|$}
        \State $(w, u, v) \gets \textsc{ExtractMin}(pq)$
        \If{$v \in visited$} \textbf{continue} \EndIf
        \State $visited.\textsc{Add}(v)$;\ $mst.\textsc{Append}(\textsc{Edge}(u,v,w))$
        \For{each $(nbr, wt)$ in \textsc{Neighbours}($G$, $v$) where $nbr \notin visited$}
            \State $pq.\textsc{Insert}(wt, v, nbr)$
        \EndFor
    \EndWhile
    \Return $mst$
\EndFunction
\end{algorithmic}
\end{algorithm}

\textbf{Complexity:} $O(|E|\log|V|)$ time, $O(|V|+|E|)$ space.

\subsubsection{Kruskal's Algorithm}

Kruskal's algorithm sorts all edges by weight and adds each edge to the
MST if its endpoints belong to different components, using the
\texttt{DisjointSet} structure for cycle detection. The loop exits early
once $|V|-1$ edges have been accepted.

\begin{algorithm}[H]
\caption{Kruskal's MST}
\begin{algorithmic}[1]
\Function{KruskalMST}{$G$}
    \State $ds \gets \textsc{DisjointSet}(V)$;\ $edges \gets \textsc{SortedByWeight}(E)$
    \State $mst \gets []$
    \For{each $e = (u, v, w)$ in $edges$}
        \If{\textsc{Union}($ds$, $u$, $v$)}
            \State $mst.\textsc{Append}(e)$
            \If{$|mst| = |V|-1$} \textbf{break} \EndIf
        \EndIf
    \EndFor
    \Return $mst$
\EndFunction
\end{algorithmic}
\end{algorithm}

\textbf{Complexity:} $O(|E|\log|E|)$ time (dominated by sorting),
$O(|V|)$ space.

\subsubsection{Comparison}

\begin{table}[H]
\centering
\caption{Prim vs.\ Kruskal.}
\begin{tabularx}{\linewidth}{@{}L{3.2cm}L{4.4cm}L{4.9cm}@{}}
\toprule
\textbf{Property} & \textbf{Prim} & \textbf{Kruskal} \\
\midrule
Approach       & Vertex-centric; grows one component & Edge-centric; merges components \\
Data structure & Min-heap                            & Disjoint-set + sorted edge list \\
Preferred for  & Dense graphs                        & Sparse graphs \\
Start node     & Required                            & Not required \\
\bottomrule
\end{tabularx}
\end{table}

% ============================================================
\section{Graph Visualiser}
\label{sec:viz}
% ============================================================

\subsection{Overview}

The visualiser (\texttt{src/visualization/pygame\_view.py}) is a
self-contained pygame application supporting two primary modes:

\begin{description}
    \item[\textbf{Edit mode}] The user constructs a graph interactively:
          left-clicking empty space creates a node; clicking two nodes in
          sequence opens a floating weight-input box to create an edge.
    \item[\textbf{Algorithm mode}] BFS, DFS, Dijkstra, Prim, or Kruskal
          is animated step-by-step with full playback controls.
\end{description}

\subsection{Interaction State Machine}

The visualiser tracks a string-valued mode with four states:

\begin{table}[H]
\centering
\caption{Visualiser mode states.}
\begin{tabular}{@{}lll@{}}
\toprule
\textbf{State} & \textbf{Trigger} & \textbf{Description} \\
\midrule
\texttt{"edit"}        & Initial / \texttt{R}                        & Graph editing enabled \\
\texttt{"pick\_start"} & \texttt{B}, \texttt{D}, \texttt{K}, \texttt{P}, \texttt{U} & Awaiting start node \\
\texttt{"pick\_end"}   & Start chosen (Dijkstra only)                & Awaiting end node \\
\texttt{"algo"}        & Node(s) confirmed                           & AlgoPlayer active \\
\bottomrule
\end{tabular}
\end{table}

All states return to \texttt{"edit"} via \texttt{Escape} or \texttt{R}.

\subsection{Animation Model}

The step generators (\texttt{\_bfs\_steps}, \texttt{\_dfs\_steps},
\texttt{\_dijkstra\_steps}, \texttt{\_prim\_steps},
\texttt{\_kruskal\_steps}) replicate each algorithm's logic exactly,
recording an \texttt{AlgoStep} snapshot after every meaningful event.
All steps are precomputed at launch, enabling instant backward stepping
and variable-speed playback without re-running the algorithm.

\begin{lstlisting}[caption={\texttt{AlgoStep} dataclass.}]
@dataclass
class AlgoStep:
    visited     : Set[str]
    frontier    : Set[str]
    active_node : Optional[str]
    active_edges: Set[Tuple[str, str]]
    active_edge : Optional[Tuple[str, str]]
    dist_labels : Dict[str, str]          # Dijkstra only
    path_nodes  : Set[str]                # Dijkstra final step
    path_edges  : Set[Tuple[str, str]]    # Dijkstra / MST final step
    mst_edges   : Set[Tuple[str, str]]    # Prim / Kruskal
    description : str
\end{lstlisting}

\subsection{Visual Encoding}

\begin{table}[H]
\centering
\caption{Node and edge colour encoding during algorithm playback.}
\begin{tabular}{@{}L{3.8cm}L{3.5cm}L{5.5cm}@{}}
\toprule
\textbf{Element} & \textbf{Colour} & \textbf{Meaning} \\
\midrule
Node -- default   & Blue (30, 120, 220)   & Unvisited \\
Node -- frontier  & Orange (220, 130, 30) & In queue / stack / heap \\
Node -- active    & White (240, 240, 255) & Being expanded this step \\
Node -- visited   & Teal (20, 160, 140)   & Fully processed \\
Node -- path      & Gold (255, 180, 20)   & On Dijkstra shortest path \\
Node -- MST       & Purple               & In MST (Prim / Kruskal) \\
Edge -- default   & Grey (140, 155, 175)  & Untraversed \\
Edge -- traversed & Teal (20, 160, 140)   & Crossed during traversal \\
Edge -- active    & White                 & Being relaxed / added now \\
Edge -- path/MST  & Gold (255, 210, 50)   & Final shortest path or MST edge \\
\bottomrule
\end{tabular}
\end{table}

For Dijkstra, the current best-known distance from the source is rendered
above each node and updated at every step.

\subsection{Playback Controls}

\begin{table}[H]
\centering
\caption{Keyboard controls during playback.}
\begin{tabular}{@{}ll@{}}
\toprule
\textbf{Key} & \textbf{Action} \\
\midrule
\texttt{Space}             & Play / Pause \\
\texttt{Right arrow}       & Step forward \\
\texttt{Left arrow}        & Step backward \\
\texttt{+} / \texttt{=}   & Increase speed \\
\texttt{-}                 & Decrease speed \\
\texttt{R} / \texttt{Esc} & Reset to edit mode \\
\bottomrule
\end{tabular}
\end{table}

Speed levels cycle through $\{0.5,\,1,\,2,\,4,\,8\}$ steps per second.
The \texttt{AlgoPlayer} time accumulator advances the step index at the
correct rate regardless of frame rate.

\subsection{Edit Mode Controls}

\begin{table}[H]
\centering
\caption{Edit mode input controls.}
\begin{tabular}{@{}ll@{}}
\toprule
\textbf{Input} & \textbf{Action} \\
\midrule
Left-click empty space          & Add node \\
Left-click node                 & Select as edge start \\
Left-click second node          & Open weight input \\
Left-click same node            & Deselect \\
Right-click node (no selection) & Delete node and all incident edges \\
Right-click edge                & Delete edge \\
Right-click (mid-selection)     & Cancel current selection \\
Scroll wheel                    & Ignored \\
\texttt{Enter}                  & Confirm edge weight \\
\texttt{Backspace}              & Delete character in weight input \\
\texttt{Escape}                 & Cancel action \\
\texttt{B}                      & Launch BFS \\
\texttt{D}                      & Launch DFS \\
\texttt{K}                      & Launch Dijkstra \\
\texttt{P}                      & Launch Prim's MST \\
\texttt{U}                      & Launch Kruskal's MST \\
\bottomrule
\end{tabular}
\end{table}

% ============================================================
\section{Random Graph Generator}
\label{sec:generator}
% ============================================================

\texttt{generate\_random\_graph} in \texttt{src/utils/graph\_generator.py}
produces a random connected weighted graph in two phases:

\begin{enumerate}
    \item \textbf{Connectivity chain} -- a Hamiltonian path
          $N_0 - N_1 - \cdots - N_{n-1}$ with random weights in
          $[1,\,\texttt{max\_weight}]$ guarantees connectivity.
    \item \textbf{Random augmentation} -- for each pair $(i,j)$ with
          $i < j$, an additional edge is added with probability
          \texttt{edge\_probability} (default 0.4).
\end{enumerate}

\begin{lstlisting}[caption={Random graph generator (\texttt{src/utils/graph\_generator.py}).}]
def generate_random_graph(n_nodes=6, edge_probability=0.4, max_weight=10):
    g     = Graph()
    nodes = [f"N{i}" for i in range(n_nodes)]
    for n in nodes:
        g.add_node(n)
    for i in range(n_nodes - 1):                   # phase 1: chain
        g.add_edge(nodes[i], nodes[i + 1], random.randint(1, max_weight))
    for i in range(n_nodes):                        # phase 2: augment
        for j in range(i + 1, n_nodes):
            if random.random() < edge_probability:
                g.add_edge(nodes[i], nodes[j], random.randint(1, max_weight))
    return g
\end{lstlisting}

The expected edge count after augmentation is
$(n-1) + \binom{n}{2}\cdot p$. For the defaults $n=6$, $p=0.4$ this
gives approximately 11 edges. Phase~1 edges may be overwritten by
Phase~2 if the same pair is selected, which is harmless since
\texttt{add\_edge} updates the weight in place.

% ============================================================
\section{Command-Line Interface}
\label{sec:cli}
% ============================================================

The \texttt{Menu} class in \texttt{src/ui/menu.py} provides a numbered
text menu that exposes the full graph API and all five algorithms:

\begin{table}[H]
\centering
\caption{CLI menu options.}
\begin{tabular}{@{}cl@{}}
\toprule
\textbf{Option} & \textbf{Action} \\
\midrule
1  & Add node \\
2  & Add edge \\
3  & Display adjacency list \\
4  & BFS traversal \\
5  & DFS traversal \\
6  & Shortest path (Dijkstra) \\
7  & Minimum spanning tree (Prim) \\
8  & Minimum spanning tree (Kruskal) \\
9  & Generate random graph \\
10 & Open pygame visualiser \\
11 & Remove node \\
12 & Remove edge \\
13 & Load sample graph \\
0  & Exit \\
\bottomrule
\end{tabular}
\end{table}

The graph object persists across menu invocations. When option 10 is
selected, the current \texttt{Graph} instance is passed directly to
\texttt{GraphVisualizer}, so any graph built or loaded via the CLI is
immediately available for visual inspection.

% ============================================================
\section{Conclusion}
\label{sec:conc}
% ============================================================

RouteForge delivers a modular, correct, and interactive graph algorithm
system. The separation of core data structures, algorithm implementations,
validation, and visualisation into independent modules means each
component can be tested, extended, or replaced in isolation.

The precomputed step-list animation model supports instant backward
stepping and variable-speed playback -- capabilities that live execution
would not allow. The consistent visual encoding across all five algorithms
(BFS, DFS, Dijkstra, Prim, Kruskal) lets users compare how each
algorithm explores the same graph under identical starting conditions.

\subsection{Potential Extensions}

\begin{enumerate}
    \item \textbf{Directed graph support} -- extend \texttt{Graph} for
          directed edges, enabling topological sort and Bellman-Ford.
    \item \textbf{A* search} -- a natural extension for geographic
          pathfinding with heuristic guidance.
    \item \textbf{Graph import/export} -- JSON or GraphML serialisation
          for persistent graphs and external tool interoperability.
    \item \textbf{Drag-and-drop node repositioning} -- mouse drag
          support for improved layout control on larger graphs.
    \item \textbf{Unit test suite} -- the \texttt{tests/} directory is
          currently empty; property-based tests using \texttt{hypothesis}
          would significantly increase correctness confidence.
\end{enumerate}

\end{document}